\input{../tex-inputs/latex-leseansicht-vorspann}

\section[Gustav Schwarzkopf an Arthur Schnitzler, 3. 10. 1899]{L04291 Gustav Schwarzkopf an Arthur Schnitzler, 3. 10. 1899}
\nopagebreak\mylabel{L04291v}
\rehead{ }\normalsize\beginnumbering\briefempfaengerindex{Schnitzler, Arthur@\textsc{Schnitzler, Arthur}!zzzSchwarzkopf, Gustav@\emph{von Gustav Schwarzkopf}!1899-10-034@{3. 10. 1899}|(be}
\toendnotes[C]{\smallbreak\pagebreak[2]}
\correspDesc{Versand  durch Gustav Schwarzkopf am 3. 10. 1899 in Wien
\newline{}Erhalt  durch Arthur Schnitzler im Zeitraum [4. 10. 1899 – 8. 10. 1899?] in Berlin}\toendnotes[C]{\smallbreak}
\Standort{CUL, Schnitzler, B 96a.}
\physDesc{Brief, 1 Blatt, 4 Seiten, 1540 Zeichen
\newline{}Handschrift: schwarze Tinte, deutsche Kurrent}\toendnotes[C]{\smallbreak}
\pstart
           \raggedleft{}{\pb}3. 10. 99\pend
           \vspace{0.5em}
\pstart
           Lieber Arthur, ich war wieder zehn Tage in der Brühl\oindex{Brühl@\textbf{Brühl}, \emph{Tal}|pw}, habe{ }ſogar von dort einen Ausflug nach Baden\oindex{Baden bei Wien@\textbf{Baden bei Wien}, \emph{Hauptstadt}|pw} und Vöslau\oindex{Bad Vöslau@\textbf{Bad Vöslau}, \emph{Hauptstadt}|pw} gemacht – nur um den So{\geminationm}er durch eine
               intereſſante Eri{\geminationn}erung zu markiren – und bin erſt heute
               nach Wien\oindex{Wien@\textbf{Wien}, \emph{Verwaltungsgebiet}|pw} zurückgeko{\geminationm}en. Habe hier Ihren lieben \label{K_L04291-1v}\edtext{Brief}{\lemma{\textnormal{\emph{Brief}}}\Cendnote{\textnormal{XXXX Auszeichnungsfehler: Dokument L04093 nicht gefunden.}}}\label{K_L04291-1} vorgefunden und gleichzeitig ihre \label{K_L04291-2v}\edtext{Karten}{\lemma{\textnormal{\emph{Karten}}}\Cendnote{\textnormal{XXXX}}}\label{K_L04291-2} aus Frankfurt\oindex{Frankfurt am Main@\textbf{Frankfurt am Main}, \emph{Hauptstadt}|pw}, und Wiesbaden\oindex{Wiesbaden@\textbf{Wiesbaden}|pw}. Die letztere {\pb}hat
               erſt den Umweg über Bukareſt\oindex{Bukarest@\textbf{Bukarest}, \emph{Verwaltungsgebiet}|pw} gemacht, wie aus dem
               Poſtstempel erſichtlich, ehe{ }ſie den Weg nach Wien\oindex{Wien@\textbf{Wien}, \emph{Verwaltungsgebiet}|pw}
               fand.\pend
           
\pstart
           – Das Verſchwinden Ihrer Stücke\pwindex{Schnitzler, Arthur 15.\,5.\,1862 Wien – 21.\,10.\,1931 ebd.@\textsc{Schnitzler, Arthur} (15.\,5.\,1862 Wien – 21.\,10.\,1931 ebd.), \emph{Schriftsteller, Mediziner}!grüne Kakadu – Paracelsus – Die Gefährtin. Drei Einakter@\strich\emph{Der grüne Kakadu – Paracelsus – Die Gefährtin. Drei Einakter}|pwv}
               habe ich natürlich bemerkt; ich muß Ihnen aufrichtig{ }ſagen, es hat mich nicht
               überraſcht. Er\pwindex{Schlenther, Paul 20.\,8.\,1854 Chernyakhovsk – 30.\,4.\,1916 Berlin@\textsc{Schlenther, Paul} (20.\,8.\,1854 Chernyakhovsk – 30.\,4.\,1916 Berlin), \emph{Schriftsteller, Kritiker, Theaterleiter}|pwv} hat Ihnen
               natürlich mit irgend einer Ausrede geantwortet und die Verſicherung gegeben, daß man
               auf ihn rechnen kö{\geminationn}e, daß er unentwegt
               u. ſ. w.\pend
           
\pstart
           Die Beka{\geminationn}ten, die Sie erwähnen, Salten\pwindex{Salten, Felix 6.\,9.\,1869 Budapest – 8.\,10.\,1945 Zürich@\textsc{Salten, Felix} (6.\,9.\,1869 Budapest – 8.\,10.\,1945 Zürich), \emph{Schriftsteller, Journalist, Chefredakteur}|pw}, Waſſerma{\geminationn}\pwindex{Wassermann, Jakob 10.\,3.\,1873 Fürth – 1.\,1.\,1934 Altaussee@\textsc{Wassermann, Jakob} (10.\,3.\,1873 Fürth – 1.\,1.\,1934 Altaussee), \emph{Schriftsteller}|pw}, u.ſ.w habe ich{ }ſeit Monaten nicht geſehen,
               auch{ }ſeit langer {\pb}Zeit kein Kaffehaus
               beſucht, keine Berliner\oindex{Berlin@\textbf{Berlin}, \emph{Hauptstadt}|pw} Zeitung geleſen, ich
               wußte{ }ſchon gar nicht mehr was vorgeht. Heute habe ich zum erſtenmal wieder Graben\oindex{Wien@\textbf{Wien}!I., Innere Stadt@\textbf{I., Innere Stadt}!Graben@\textbf{Graben}, \emph{Straße}|pw}luft geatmet und mit einigen Sta{\geminationm}gäſten geſprochen.\pend
           
\pstart
           Ich \uline{wünſche} es gewiß, daß Sie um den 10.
               herum nach Wien\oindex{Wien@\textbf{Wien}, \emph{Verwaltungsgebiet}|pw} ko{\geminationm}en, aber ich glaube noch nicht recht daran. Man wird verſuchen Sie länger in Berlin\oindex{Berlin@\textbf{Berlin}, \emph{Hauptstadt}|pw} feſtzuhalten und da unter »man« die
               verſchiedenſten weiblichen Wesen\pwindex{Elsinger, Marie *~28.\,2.\,1874 St. Pölten@\textsc{Elsinger, Marie} (*~28.\,2.\,1874 St. Pölten), \emph{Schauspielerin}|pwv} zu verſtehen{ }ſind,{ }ſo werde ich wol leider Recht behalten.\pend
           
\pstart
           {\pb}Über \textsc{Dreyfus\pwindex{Dreyfus, Alfred 9.\,10.\,1859 Mulhouse – 12.\,7.\,1935 Paris@\textsc{Dreyfus, Alfred} (9.\,10.\,1859 Mulhouse – 12.\,7.\,1935 Paris), \emph{Militär}|pw}}, und \textsc{Polna\oindex{Polná@\textbf{Polná}|pw}}, über den kleinen Kraus\pwindex{Kraus, Karl 28.\,4.\,1874 Jičín – 12.\,6.\,1936 Wien@\textsc{Kraus, Karl} (28.\,4.\,1874 Jičín – 12.\,6.\,1936 Wien), \emph{Schriftsteller, Publizist, Schriftsteller}|pw}, die Fackel\pwindex{Fackel@\emph{Die Fackel}|pw} und die i{\geminationm}er
               mehr zunehmende Wiener\oindex{Wien@\textbf{Wien}, \emph{Verwaltungsgebiet}|pw} Vertrottelung{ }ſowie über
               die vielen anderen Dinge, die \textcolor{gray}{Anlaß} zu Ärger und Ekel geben,{ }ſprechen wir, we{\geminationn} Sie einmal da{ }ſind.\pend
           
\pstart
           Bis dahin{\\[\baselineskip]}herzlichſt{\\[\baselineskip]}Ihr{\\[\baselineskip]}\spacefill\mbox{Gustav}\pend
           \leftskip=0em{}
\pstart
           \noindent{}Den Damen Glümer\pwindex{Glümer, Marie 3.\,7.\,1867 Wien – 16.\,11.\,1925 München@\textsc{Glümer, Marie} (3.\,7.\,1867 Wien – 16.\,11.\,1925 München), \emph{Schauspielerin}|pw}\pwindex{Glümer, Auguste 16.\,3.\,1862 Wien – 1956@\textsc{Glümer, Auguste} (16.\,3.\,1862 Wien – 1956), \emph{Lehrerin}|pw} meine
                  Empfehlungen.\pend
           \selectlanguage{ngerman}\endnumbering\briefempfaengerindex{Schnitzler, Arthur@\textsc{Schnitzler, Arthur}!zzzSchwarzkopf, Gustav@\emph{von Gustav Schwarzkopf}!1899-10-034@{3. 10. 1899}|)be}\mylabel{L04291h}
\begin{anhang}
\end{anhang}\newcommand{\dateiname}{L04291}\newcommand{\titel}{Gustav Schwarzkopf an Arthur Schnitzler, 3. 10. 1899}\newcommand{\editorInnen}{Herausgegeben von Jahnke, SelmaMüller, Martin Anton}\input{../tex-inputs/latex-leseansicht-abspann}
